\documentclass[12pt,a4paper]{article}
\usepackage[spanish]{babel}
\usepackage[utf8]{inputenc}
\usepackage[T1]{fontenc}
\usepackage{lmodern}
\usepackage[table]{xcolor}
\usepackage{geometry}
\usepackage{setspace}
\usepackage{longtable}
\usepackage{booktabs}
\usepackage{array}
\usepackage{tabularx}
\usepackage{hyperref}
\usepackage{enumitem}

\geometry{margin=2.4cm}
\setstretch{1.12}
\setlength{\parindent}{0pt}
\setlength{\parskip}{0.55em}
\hypersetup{colorlinks=true, linkcolor=black, urlcolor=blue, citecolor=black}

\definecolor{PrimaryBlue}{HTML}{16324F}
\definecolor{SoftGray}{HTML}{F3F6F8}

	itle{\Huge\textbf{Modelo Entidad-Relaci\'on}\[0.3em]\Large Plataforma Web de Mascotas 3D}
\author{Equipo del Proyecto}
\date{Mayo 2026}

\begin{document}

\begin{titlepage}
\centering
\vspace*{1.8cm}
{\Huge\bfseries\color{PrimaryBlue} MER\\[0.35em]}
{\LARGE Modelo Entidad-Relaci\'on y dise\~no de Base de Datos\\[1.2cm]}

\begin{tabular}{p{5cm} p{8cm}}
\rowcolor{SoftGray}
\textbf{Versi\'on} & 1.0 \\
\textbf{Fecha} & Mayo 2026 \\
\textbf{Motor BD} & PostgreSQL \\
\textbf{Alcance} & MVP \\
\end{tabular}

\vfill
{\large Documento para compilaci\'on en Overleaf}
\end{titlepage}

\tableofcontents
\newpage

\section{Introducci\'on}
Este documento especifica el modelo de datos de la plataforma. Incluye entidades, atributos, relaciones y restricciones.

\section{Diagrama conceptual}

A nivel conceptual, el modelo incluye:
\begin{itemize}[leftmargin=*]
    \item \textbf{Usuarios}: Personas que se registran en la plataforma.
    \item \textbf{Mascotas}: Registros de mascotas propiedad de un usuario.
    \item \textbf{Modelos 3D}: Archivos de modelos base asociados a tipos de animales.
    \item \textbf{Fotos}: Im\'agenes de mascotas subidas por usuarios.
\end{itemize}

\section{Entidades y atributos}

\subsection{Entidad: Usuarios}

\begin{table}[htbp]
\centering
\renewcommand{\arraystretch}{1.1}
\begin{tabularx}{\textwidth}{>{\bfseries}p{2.5cm} p{1.5cm} p{1.5cm} X}
\toprule
Atributo & Tipo & Clave & Descripci\'on \\
\midrule
id & UUID & PK & Identificador \'unico. \\
email & VARCHAR(255) & UNIQUE & Correo electr\'onico \'unico. \\
password\_hash & VARCHAR(255) & & Contrase\~na hasheada (Bcrypt). \\
nombre & VARCHAR(100) & & Nombre del usuario (opcional). \\
created\_at & TIMESTAMP & & Fecha y hora de creaci\'on. \\
updated\_at & TIMESTAMP & & Fecha y hora de \'ultima modificaci\'on. \\
\bottomrule
\end{tabularx}
\end{table}

\subsection{Entidad: Mascotas}

\begin{table}[htbp]
\centering
\renewcommand{\arraystretch}{1.1}
\begin{tabularx}{\textwidth}{>{\bfseries}p{2.5cm} p{1.5cm} p{1.5cm} X}
\toprule
Atributo & Tipo & Clave & Descripci\'on \\
\midrule
id & UUID & PK & Identificador \'unico. \\
usuario\_id & UUID & FK & Referencia a Usuarios. \\
nombre & VARCHAR(100) & & Nombre de la mascota. \\
edad & INTEGER & & Edad en a\~nos. \\
tipo\_animal & VARCHAR(50) & & Gato, Perro, Jaguar, \'Aguila, Serpiente, etc. \\
descripcion & TEXT & & Informaci\'on adicional. \\
created\_at & TIMESTAMP & & Fecha de creaci\'on. \\
updated\_at & TIMESTAMP & & Fecha de \'ultima modificaci\'on. \\
\bottomrule
\end{tabularx}
\end{table}

\subsection{Entidad: Modelos 3D}

\begin{table}[htbp]
\centering
\renewcommand{\arraystretch}{1.1}
\begin{tabularx}{\textwidth}{>{\bfseries}p{2.5cm} p{1.5cm} p{1.5cm} X}
\toprule
Atributo & Tipo & Clave & Descripci\'on \\
\midrule
id & UUID & PK & Identificador \'unico. \\
tipo\_animal & VARCHAR(50) & & Tipo de animal (Gato, Perro, etc.). \\
nombre & VARCHAR(100) & & Nombre descriptivo. \\
archivo\_nombre & VARCHAR(255) & & Nombre del archivo GLB/GLTF. \\
ruta\_almacenamiento & VARCHAR(500) & & Ruta relativa en servidor. \\
tama\~no\_bytes & BIGINT & & Tama\~no del archivo. \\
mime\_type & VARCHAR(50) & & Tipo MIME (model/gltf-binary). \\
created\_at & TIMESTAMP & & Fecha de creaci\'on. \\
\bottomrule
\end{tabularx}
\end{table}

\subsection{Entidad: Fotos}

\begin{table}[htbp]
\centering
\renewcommand{\arraystretch}{1.1}
\begin{tabularx}{\textwidth}{>{\bfseries}p{2.5cm} p{1.5cm} p{1.5cm} X}
\toprule
Atributo & Tipo & Clave & Descripci\'on \\
\midrule
id & UUID & PK & Identificador \'unico. \\
mascota\_id & UUID & FK & Referencia a Mascotas. \\
nombre\_archivo & VARCHAR(255) & & Nombre del archivo. \\
ruta\_almacenamiento & VARCHAR(500) & & Ruta en servidor. \\
tama\~no\_bytes & BIGINT & & Tama\~no de la imagen. \\
mime\_type & VARCHAR(50) & & Tipo MIME (image/jpeg, image/png). \\
created\_at & TIMESTAMP & & Fecha de carga. \\
\bottomrule
\end{tabularx}
\end{table}

\section{Relaciones}

\subsection{Relaci\'on: Usuario -- Mascota}
\begin{itemize}[leftmargin=*]
    \item \textbf{Tipo}: Uno a Muchos (1:N).
    \item \textbf{Descripci\'on}: Un usuario puede tener m\'ultiples mascotas. Una mascota pertenece a un solo usuario.
    \item \textbf{Implementaci\'on}: mascota.usuario\_id $\rightarrow$ usuario.id (FK).
    \item \textbf{Integridad}: ON DELETE CASCADE (si usuario se elimina, sus mascotas se eliminan).
\end{itemize}

\subsection{Relaci\'on: Mascota -- Foto}
\begin{itemize}[leftmargin=*]
    \item \textbf{Tipo}: Uno a Muchos (1:N).
    \item \textbf{Descripci\'on}: Una mascota puede tener m\'ultiples fotos. Una foto pertenece a una sola mascota.
    \item \textbf{Implementaci\'on}: foto.mascota\_id $\rightarrow$ mascota.id (FK).
    \item \textbf{Integridad}: ON DELETE CASCADE.
\end{itemize}

\subsection{Relaci\'on: Mascota -- Modelo 3D}
\begin{itemize}[leftmargin=*]
    \item \textbf{Tipo}: Muchos a Uno (N:1).
    \item \textbf{Descripci\'on}: M\'ultiples mascotas pueden asociarse al mismo modelo 3D base seg\'un su tipo de animal.
    \item \textbf{Implementaci\'on}: mascota.tipo\_animal $\rightarrow$ modelo3d.tipo\_animal (FK l\'ogica).
\end{itemize}

\section{Trazabilidad funcional}
\begin{itemize}[leftmargin=*]
    \item RF-01 y RF-02: Entidad Usuarios.
    \item RF-03 y RF-04: Entidad Mascotas.
    \item RF-05: Entidad Mascotas, Fotos y Modelos 3D.
    \item RF-06: Entidad Fotos.
    \item RF-07: Entidad Modelos 3D y l\'ogica de administraci\'on.
\end{itemize}

\subsection{Observaci\'on sobre sesi\'on}
La expiraci\'on por inactividad (RNF-05) se gestiona en la capa de aplicaci\'on y no requiere una entidad persistente dedicada en el modelo conceptual.

\section{Restricciones}

\subsection{Restricciones de dominio}
\begin{itemize}[leftmargin=*]
    \item email en Usuarios: debe ser \'unico y v\'alido.
    \item edad en Mascotas: debe ser $\geq 0$.
    \item tipo\_animal: solo valores permitidos (Gato, Perro, Jaguar, \'Aguila, Serpiente).
    \item tama\~no\_bytes en Fotos: $\leq$ 5242880 (5MB).
\end{itemize}

\subsection{Restricciones de integridad}
\begin{itemize}[leftmargin=*]
    \item Integridad referencial: mascota.usuario\_id debe existir en usuarios.id.
    \item Integridad referencial: foto.mascota\_id debe existir en mascotas.id.
    \item Unicidad: (usuario.email) \'unico.
\end{itemize}

\section{Script SQL de creaci\'on}

\subsection{Tabla Usuarios}
\begin{verbatim}
CREATE TABLE usuarios (
  id UUID PRIMARY KEY DEFAULT gen_random_uuid(),
  email VARCHAR(255) UNIQUE NOT NULL,
  password_hash VARCHAR(255) NOT NULL,
  nombre VARCHAR(100),
  created_at TIMESTAMP DEFAULT CURRENT_TIMESTAMP,
  updated_at TIMESTAMP DEFAULT CURRENT_TIMESTAMP
);
\end{verbatim}

\subsection{Tabla Mascotas}
\begin{verbatim}
CREATE TABLE mascotas (
  id UUID PRIMARY KEY DEFAULT gen_random_uuid(),
  usuario_id UUID NOT NULL REFERENCES usuarios(id) ON DELETE CASCADE,
  nombre VARCHAR(100) NOT NULL,
  edad INTEGER,
  tipo_animal VARCHAR(50) NOT NULL,
  descripcion TEXT,
  created_at TIMESTAMP DEFAULT CURRENT_TIMESTAMP,
  updated_at TIMESTAMP DEFAULT CURRENT_TIMESTAMP
);

CREATE INDEX idx_mascotas_usuario_id ON mascotas(usuario_id);
\end{verbatim}

\subsection{Tabla Modelos 3D}
\begin{verbatim}
CREATE TABLE modelos3d (
  id UUID PRIMARY KEY DEFAULT gen_random_uuid(),
  tipo_animal VARCHAR(50) UNIQUE NOT NULL,
  nombre VARCHAR(100) NOT NULL,
  archivo_nombre VARCHAR(255) NOT NULL,
  ruta_almacenamiento VARCHAR(500) NOT NULL,
  tamano_bytes BIGINT,
  mime_type VARCHAR(50),
  created_at TIMESTAMP DEFAULT CURRENT_TIMESTAMP
);
\end{verbatim}

\subsection{Tabla Fotos}
\begin{verbatim}
CREATE TABLE fotos (
  id UUID PRIMARY KEY DEFAULT gen_random_uuid(),
  mascota_id UUID NOT NULL REFERENCES mascotas(id) ON DELETE CASCADE,
  nombre_archivo VARCHAR(255) NOT NULL,
  ruta_almacenamiento VARCHAR(500) NOT NULL,
  tamano_bytes BIGINT,
  mime_type VARCHAR(50),
  created_at TIMESTAMP DEFAULT CURRENT_TIMESTAMP
);

CREATE INDEX idx_fotos_mascota_id ON fotos(mascota_id);
\end{verbatim}

\section{Consideraciones de dise\~no}

\subsection{Normalizaci\'on}
\begin{itemize}[leftmargin=*]
    \item El modelo est\'a en Tercera Forma Normal (3NF).
    \item No hay redundancia innecesaria.
    \item Relaciones bien definidas.
\end{itemize}

\subsection{\'Indices}
\begin{itemize}[leftmargin=*]
    \item \'Indice en usuarios.email para b\'usquedas de login r\'apidas.
    \item \'Indice en mascotas.usuario\_id para filtrados por usuario.
    \item \'Indice en fotos.mascota\_id para relaciones.
\end{itemize}

\subsection{Escalabilidad}
\begin{itemize}[leftmargin=*]
    \item UUIDs en lugar de enteros para distribuci\'on futura.
    \item Timestamps para auditor\'ia y control de cambios.
    \item Dise\~no sin dependen cias circulares.
\end{itemize}

\end{document}
